\begin{proofbox}
\textbf{证明一（优先）：代入、展开、韦达。}
先证横开口情形。若 $B\ne0$，由直线得到 $y=-(Ax+C)/B$，所以
\[
 \frac{x^2}{a^2}-\frac{(Ax+C)^2}{b^2B^2}=1
 \ \Longleftrightarrow\
 b^2B^2x^2-a^2(Ax+C)^2=a^2b^2B^2.
\]
展开并合并得到
\[
 Mx^2-2ACa^2x-a^2(C^2+B^2b^2)=0.
\]
若 $M\ne0$，由韦达定理直接得到 $x_1+x_2$ 和 $x_1x_2$。
这是等价变形，每个实根都唯一恢复出 $y=-(Ax+C)/B$，无增根或漏点。
若忘记韦达，只要展开 $p(x-x_1)(x-x_2)$，比较一次项及常数项即可。

若 $A\ne0$，改用 $x=-(By+C)/A$，代入后得到
\[
 b^2(By+C)^2-a^2A^2y^2=a^2b^2A^2,
\]
展开为 $My^2+2BCb^2y+b^2(C^2-A^2a^2)=0$，再用韦达。

\textbf{交点分类也只需化简一次。} 在 $B\ne0$ 时，
\[
 \begin{aligned}
 \Delta_x
 &=4A^2C^2a^4+4Ma^2(C^2+B^2b^2)\\
 &=4a^2\bigl[C^2(A^2a^2+M)+MB^2b^2\bigr]\\
 &=4a^2b^2B^2(C^2+M)=4a^2b^2B^2H.
 \end{aligned}
\]
在 $A\ne0$ 时，同理 $\Delta_y=4a^2b^2A^2H$。
选择合法消元方向后，前面的因子为正，故 $H$ 的符号恰好给出三类实交点情况。
当 $M=0$ 时，$b^2B^2=a^2A^2$ 且 $A,B\ne0$：$C\ne0$ 时一次项不为零，\textbf{恰有}一个实根；$C=0$ 时常数项非零，无解。
\end{proofbox}
\clearpage
\begin{proofbox}
\textbf{证明二（少记公式）：交换坐标，一组推出另一组。}
令 $X=y,Y=x$。纵开口的双曲线变为横开口：
\[
 \frac{X^2}{a^2}-\frac{Y^2}{b^2}=1,\qquad BX+AY+C=0.
\]
\begin{center}\begin{tikzpicture}[font=\small]
\draw[black!30,rounded corners=2pt] (0,0) rectangle (14.4,8.2);
\node[anchor=west,font=\small\bfseries] at (.5,7.72) {交换坐标：第二组公式无需重新背一遍};
\begin{scope}[shift={(3.7,3.7)},scale=0.57]
\begin{scope}
\clip (-3.6,-3.6) rectangle (3.6,3.6);
\draw[blue,very thick,domain=-3.6:3.6,samples=161,variable=\t] plot ({2*sqrt(1+(\t/1)^2)},\t);
\draw[blue,very thick,domain=-3.6:3.6,samples=161,variable=\t] plot ({-2*sqrt(1+(\t/1)^2)},\t);
\end{scope}
\draw[-Stealth,black!45] (-3.6,0)--(3.6,0) node[right] {$x$};
\draw[-Stealth,black!45] (0,-3.6)--(0,3.6) node[above] {$y$};
\draw[orange!85!black,very thick] (-3.6,-1.6)--(.1,2.1);\coordinate (P) at (-3.3333333333,-1.3333333333);\fill[black] (P) circle (1.8pt) node[above left=6pt] {$P$};
\coordinate (Q) at (-2.0000000000,0.0000000000);\fill[black] (Q) circle (1.8pt) node[below right=9pt] {$Q$};
\end{scope}\begin{scope}[shift={(10.6,3.7)},scale=0.57]
\begin{scope}
\clip (-3.6,-3.6) rectangle (3.6,3.6);
\draw[blue,very thick,domain=-3.6:3.6,samples=161,variable=\t] plot (\t,{2*sqrt(1+(\t/1)^2)});
\draw[blue,very thick,domain=-3.6:3.6,samples=161,variable=\t] plot (\t,{-2*sqrt(1+(\t/1)^2)});
\end{scope}
\draw[-Stealth,black!45] (-3.6,0)--(3.6,0) node[right] {$x$};
\draw[-Stealth,black!45] (0,-3.6)--(0,3.6) node[above] {$y$};
\draw[orange!85!black,very thick] (-1.6,-3.6)--(2.1,.1);\coordinate (P') at (-1.3333333333,-3.3333333333);\fill[black] (P') circle (1.8pt) node[above left=10pt] {$P'$};
\coordinate (Q') at (0.0000000000,-2.0000000000);\fill[black] (Q') circle (1.8pt) node[above left=11pt] {$Q'$};
\end{scope}\node at (3.7,6.7) {$x^2/4-y^2=1$};\node at (10.6,6.7) {$y^2/4-x^2=1$};
\node[anchor=west] at (.7,.55) {左：$y=x+2$；右：$x=y+2$；交点 $(x_i,y_i)$ 对应 $(y_i,x_i)$};
\end{tikzpicture}\end{center}

在已证公式中，同时作 $x\leftrightarrow y$、$A\leftrightarrow B$。
原来的 $M$ 变成 $A^2b^2-B^2a^2=N$，$H$ 变成 $K=C^2+N$；横坐标和、积变成纵坐标和、积。
这就逐项得到纵开口的两条消元方程及四个韦达关系。
相应判别式仍为 $4a^2b^2B^2K$（消去 $y$）或 $4a^2b^2A^2K$（消去 $x$），故分类及退化结论也全部成立。
\end{proofbox}
\clearpage
\begin{proofbox}
\textbf{证明三（数形结合）：渐近线因式分解，直接看降次。}
横开口双曲线可写为
\[
 \left(\frac{x}{a}-\frac{y}{b}\right)
 \left(\frac{x}{a}+\frac{y}{b}\right)=1.
\]
记 $u=x/a-y/b$、$v=x/a+y/b$，则 $uv=1$。
两条渐近线分别是 $u=0$、$v=0$。
$M=0$ 意味着 $A/B=\pm b/a$，所以直线正好固定其中一个量为常数。
若固定 $u=d\ne0$，则 $v=1/d$ 被唯一确定，进而
\[
 x=\frac a2\left(d+\frac1d\right),\qquad
 y=\frac b2\left(\frac1d-d\right).
\]
因此恰有一个交点；若 $d=0$，则 $uv=1$ 无解。固定 $v$ 的情形相同。
这解释了前面渐近线平行图中的现象，不需要解二次方程。

\textbf{中点另有一条短证：两点方程作差。}
设 $M\ne0,H>0$，弦中点为 $S(s,t)$。两点的双曲线方程相减：
\[
 \frac{(x_1-x_2)(x_1+x_2)}{a^2}
 -\frac{(y_1-y_2)(y_1+y_2)}{b^2}=0.
\]
两点差向量沿直线方向 $(B,-A)$，所以
\[
 \frac{Bs}{a^2}+\frac{At}{b^2}=0,\qquad As+Bt+C=0.
\]
不需要除以 $A$ 或 $B$，解这个二元一次方程组就得到
\[
 S=\left(\frac{ACa^2}{M},-\frac{BCb^2}{M}\right).
\]
这与韦达求平均的结果一致，且覆盖水平、竖直弦；$H=0$ 时直接代入可验证得到切点。
纵开口情形由交换坐标得到。
\end{proofbox}
\clearpage
\begin{proofbox}
\textbf{弦长推导：把根差补成真正的距离。}
当 $B\ne0$，$y_2-y_1=-(A/B)(x_2-x_1)$，因此
\[
 PQ=\frac{\sqrt{A^2+B^2}}{|B|}\,|x_2-x_1|.
\]
\begin{center}\begin{tikzpicture}[font=\small]
\draw[black!30,rounded corners=2pt] (0,0) rectangle (14.4,8.2);
\node[anchor=west,font=\small\bfseries] at (.5,7.72) {横坐标根差只是投影：弦长由直角三角形得到};
\begin{scope}[shift={(7.2,3.65)},scale=1.3]
\begin{scope}
\clip (-4.5,-1.9) rectangle (4.5,1.9);
\draw[blue,very thick,domain=-1.9:1.9,samples=161,variable=\t] plot ({2*sqrt(1+(\t/1)^2)},\t);
\draw[blue,very thick,domain=-1.9:1.9,samples=161,variable=\t] plot ({-2*sqrt(1+(\t/1)^2)},\t);
\end{scope}
\draw[-Stealth,black!45] (-4.5,0)--(4.5,0) node[right] {$x$};
\draw[-Stealth,black!45] (0,-1.9)--(0,1.9) node[above] {$y$};
\coordinate (P) at (-3.3333333333,-1.3333333333);\fill[black] (P) circle (1.8pt) node[above left=10pt] {$P$};
\coordinate (Q) at (-2.0000000000,0.0000000000);\fill[black] (Q) circle (1.8pt) node[above right=7pt] {$Q$};
\coordinate (T) at (-2.0000000000,-1.3333333333);\fill[black] (T) circle (1.8pt) node[below right=7pt] {$T$};
\draw[orange!85!black,very thick] (P)--(Q);
\draw[teal!70!black,dashed,thick] (P)--(T)--(Q);
\draw (-2.16,-1.3333333333)--(-2.16,-1.1733333333)--(-2,-1.1733333333);
\node[below=11pt] at (-2.6666666667,-1.3333333333) {$|x_2-x_1|$};
\node[anchor=west] at (-1.5,-.65) {$|y_2-y_1|$};
\end{scope}
\node[anchor=west] at (.7,.55) {$PQ^2=(x_2-x_1)^2+(y_2-y_1)^2$；本例两条直角边都长 $4/3$};
\end{tikzpicture}\end{center}

由韦达或求根公式，$|x_2-x_1|=\sqrt{\Delta_x}/|M|$。
代入 $\Delta_x=4a^2b^2B^2H$，得
\[
 PQ=\frac{2ab\sqrt{(A^2+B^2)H}}{|M|}.
\]
$B=0$ 时改用 $y$ 坐标及 $\Delta_y$，得到同一公式；纵开口再交换坐标即可。
这里绝对值不可去掉，因为 $M$ 可以是负数。
\end{proofbox}
\clearpage
% ===== 新增证明四：两点相减求横坐标和（原证明一至三及 TikZ 均不变）=====
\begin{proofbox}
\textbf{证明四（考场速推）：两点相减，不展开消元二次式，直接求 $x_1+x_2$。}
考虑
\[
 \frac{x^2}{a^2}-\frac{y^2}{b^2}=1,\qquad y=kx+m.
\]
设两个\textbf{不同}交点为 $P(x_1,y_1),Q(x_2,y_2)$。
由于直线不是竖直线，两点不同就有 $x_1\ne x_2$。
两点各自在双曲线上：
\[
 \frac{x_1^2}{a^2}-\frac{y_1^2}{b^2}=1\quad(1),\qquad
 \frac{x_2^2}{a^2}-\frac{y_2^2}{b^2}=1\quad(2).
\]
\textbf{第一步：$(1)-(2)$，用平方差。}
\[
 \frac{(x_1-x_2)(x_1+x_2)}{a^2}
 =\frac{(y_1-y_2)(y_1+y_2)}{b^2}.
\]
\textbf{第二步：直线的两点纵坐标差为 $y_1-y_2=k(x_1-x_2)$。}
约去非零的 $x_1-x_2$，立即得到
\[
 \boxed{\frac{x_1+x_2}{a^2}=\frac{k(y_1+y_2)}{b^2}}.
\]
\textbf{第三步：两点的直线方程相加，$y_1+y_2=k(x_1+x_2)+2m$。}
代入上式：
\[
 \left(\frac1{a^2}-\frac{k^2}{b^2}\right)(x_1+x_2)
 =\frac{2km}{b^2}.
\]
只要 $b^2-a^2k^2\ne0$，便有
\[
 \boxed{x_1+x_2=\frac{2a^2km}{b^2-a^2k^2}}.
\]
若原直线写为 $Ax+By+C=0$（$B\ne0$），
令 $k=-A/B$、$m=-C/B$，并记 $M=B^2b^2-A^2a^2$，即恢复
\[
 \boxed{x_1+x_2=\frac{2ACa^2}{M}}.
\]
\textbf{本法的低门槛所在：} 平方差带出\emph{和}，直线关系把 $y$ 的和换成 $x$ 的和；
全程不需要完整联立、求根，也不必先写韦达式。
\end{proofbox}
\clearpage

% ===== 新增证明五：交叉乘减求横坐标积 =====
\begin{proofbox}
\textbf{证明五（考场理解法）：交叉乘减，消一次项，直接求 $x_1x_2$。}
沿用证明四的双曲线、直线和两个不同交点。由
\[
 F(x)=\frac{x^2}{a^2}-\frac{(kx+m)^2}{b^2}-1,
 \qquad F(x_1)=F(x_2)=0,
\]
\textbf{第一个方程乘 $x_2$，第二个方程乘 $x_1$，然后相减：}
\[
 \begin{aligned}
 0={}&x_2\left[\frac{x_1^2}{a^2}-\frac{(kx_1+m)^2}{b^2}-1\right]\\
 &-x_1\left[\frac{x_2^2}{a^2}-\frac{(kx_2+m)^2}{b^2}-1\right].
 \end{aligned}
\]
\textbf{关键只看平方内的交叉项：}
$x_2(2kmx_1)=x_1(2kmx_2)=2kmx_1x_2$，
所以相减时它们\textbf{完全抵消}。其余部分为
\[
 \begin{aligned}
 &x_2(kx_1+m)^2-x_1(kx_2+m)^2\\
 &\qquad=k^2x_1x_2(x_1-x_2)+m^2(x_2-x_1).
 \end{aligned}
\]
代回并提取公因子 $x_1-x_2$：
\[
 0=(x_1-x_2)\left[
 \left(\frac1{a^2}-\frac{k^2}{b^2}\right)x_1x_2
 +1+\frac{m^2}{b^2}\right].
\]
因为 $x_1\ne x_2$，可约去 $x_1-x_2$：
\[
 \boxed{\left(\frac1{a^2}-\frac{k^2}{b^2}\right)x_1x_2
 =-1-\frac{m^2}{b^2}}.
\]
若 $b^2-a^2k^2\ne0$，则
\[
 \boxed{x_1x_2=-\frac{a^2(b^2+m^2)}{b^2-a^2k^2}}.
\]
将 $k=-A/B,m=-C/B$（$B\ne0$）代入，得到原结论
\[
 \boxed{x_1x_2=-\frac{a^2(B^2b^2+C^2)}{M}},\qquad
 M=B^2b^2-A^2a^2.
\]
\textbf{为什么想到这一招？} 两点方程直接减可消常数、留下\emph{和}；
交叉乘减可消一次项、留下\emph{积}。这也是韦达关系的一种原理性推导。
\end{proofbox}
\clearpage

% ===== 再补充：证明四、五逐行不跳步版（此前全部证明及 TikZ 原样保留）=====
\begin{proofbox}
\textbf{证明四补充（每一步写全）：从“两点相减”得到横坐标和。}

\textbf{已知：} $x^2/a^2-y^2/b^2=1$，直线 $y=kx+m$；
$P(x_1,y_1)$、$Q(x_2,y_2)$ 为两个不同实交点。

\textbf{第 1 步：} 分别代入双曲线，列出两个等式：
\[
\begin{aligned}
 (1)\quad &\frac{x_1^2}{a^2}-\frac{y_1^2}{b^2}=1,\\
 (2)\quad &\frac{x_2^2}{a^2}-\frac{y_2^2}{b^2}=1.
\end{aligned}
\]
\textbf{第 2 步：} $(1)-(2)$，先让右侧常数消失：
\[
 \frac{x_1^2-x_2^2}{a^2}-\frac{y_1^2-y_2^2}{b^2}=1-1=0.
\]
\textbf{第 3 步：} 将两个平方差分别分解：
\[
 \frac{(x_1-x_2)(x_1+x_2)}{a^2}
 =\frac{(y_1-y_2)(y_1+y_2)}{b^2}.
\]
\textbf{第 4 步：} 两点分别满足直线方程：
\[
 y_1=kx_1+m,\qquad y_2=kx_2+m.
\]
这两个式子相减：
\[
 y_1-y_2=(kx_1+m)-(kx_2+m)=k(x_1-x_2).
\]
代回第 3 步：
\[
 \frac{(x_1-x_2)(x_1+x_2)}{a^2}
 =\frac{k(x_1-x_2)(y_1+y_2)}{b^2}.
\]
\textbf{第 5 步：} 非竖直直线上的两个不同点有 $x_1\ne x_2$，
故两边同除以 $x_1-x_2$，得到
\[
 \boxed{\frac{x_1+x_2}{a^2}=\frac{k(y_1+y_2)}{b^2}}.
\]
\textbf{第 6 步：} 再把直线方程\emph{相加}：
\[
 y_1+y_2=(kx_1+m)+(kx_2+m)=k(x_1+x_2)+2m.
\]
代回第 5 步：
\[
 \frac{x_1+x_2}{a^2}
 =\frac{k\,[k(x_1+x_2)+2m]}{b^2}.
\]
移项得到
\[
 \left(\frac1{a^2}-\frac{k^2}{b^2}\right)(x_1+x_2)
 =\frac{2km}{b^2}.
\]
当 $b^2-a^2k^2\ne0$ 时，解得
\[
 \boxed{x_1+x_2=\frac{2a^2km}{b^2-a^2k^2}}.
\]
\textbf{速写还原：} 曲线差 $\to$ 平方差 $\to$ 直线差约分
$\to$ 直线和代入。\textbf{不需要先把两个方程联立展开。}
\end{proofbox}
\clearpage

\begin{proofbox}
\textbf{证明五补充（每一步写全）：为什么交叉乘减产生乘积。}

\textbf{第 1 步：} 将直线 $y=kx+m$ 代入双曲线。
为方便书写，把等号右端移到左边：
\[
 \frac{x^2}{a^2}-\frac{(kx+m)^2}{b^2}-1=0.
\]
两个不同交点满足
\[
\begin{aligned}
 (1)\quad &\frac{x_1^2}{a^2}-\frac{(kx_1+m)^2}{b^2}-1=0,\\
 (2)\quad &\frac{x_2^2}{a^2}-\frac{(kx_2+m)^2}{b^2}-1=0.
\end{aligned}
\]
\textbf{第 2 步：} 把平方先\emph{部分}打开，只看出一次项的形状：
\[
 (kx_i+m)^2=k^2x_i^2+2kmx_i+m^2\quad(i=1,2).
\]
我们想消掉 $2kmx_i$，因此故意让 $(1)$ 乘 $x_2$、$(2)$ 乘 $x_1$。
下面直接把\textbf{两个新等式完整写出}：
\[
\begin{aligned}
 (1)\times x_2:\quad
 &\frac{x_1^2x_2}{a^2}
 -\frac{k^2x_1^2x_2+2kmx_1x_2+m^2x_2}{b^2}
 -x_2=0,\\
 (2)\times x_1:\quad
 &\frac{x_2^2x_1}{a^2}
 -\frac{k^2x_2^2x_1+2kmx_2x_1+m^2x_1}{b^2}
 -x_1=0.
\end{aligned}
\]
\textbf{第 3 步：} 上式第一行减第二行。请特别观察三组：
\[
\begin{aligned}
 &\text{平方项： }x_1^2x_2-x_2^2x_1=x_1x_2(x_1-x_2),\\
 &\text{一次项： }2kmx_1x_2-2kmx_2x_1=0,\\
 &\text{常数项： }-x_2-(-x_1)=x_1-x_2.
\end{aligned}
\]
$-m^2x_i/b^2$ 的交叉乘部分为
$-m^2(x_2-x_1)/b^2=+m^2(x_1-x_2)/b^2$。
于是整式变成
\[
\begin{aligned}
0={}&\frac{x_1x_2(x_1-x_2)}{a^2}
 -\frac{k^2x_1x_2(x_1-x_2)}{b^2}\\
 &+\frac{m^2(x_1-x_2)}{b^2}+(x_1-x_2).
\end{aligned}
\]
\textbf{第 4 步：} 四项\emph{全部}都有公因子 $x_1-x_2$：
\[
0=(x_1-x_2)\left[
 \left(\frac1{a^2}-\frac{k^2}{b^2}\right)x_1x_2
 +\frac{m^2}{b^2}+1\right].
\]
不同交点且直线非竖直，故 $x_1-x_2\ne0$；
去掉这一因子，剩下
\[
\boxed{\left(\frac1{a^2}-\frac{k^2}{b^2}\right)x_1x_2
=-1-\frac{m^2}{b^2}}.
\]
若 $b^2-a^2k^2\ne0$，再解出
\[
\boxed{x_1x_2=-\frac{a^2(b^2+m^2)}{b^2-a^2k^2}}.
\]
\textbf{动机再说一遍：} 求和时相减，是为了产生\emph{平方差}；
求积时交叉乘减，是为了使\emph{一次项相等后抵消}。
两者不是两个需要硬背的复杂公式，而是\textbf{两种有目标的消元动作}。
\end{proofbox}
\clearpage
